% GAL1897_W10 — source-faithful transcription of physical PDF pages 080–091.
% Primary authority: JP2 leaves L0079–L0090, navigated against the canonical
% 96-page 1897 PDF and inspected with immediate-neighbor pages PDF079/PDF092.
% Printed and substantive defects remain unchanged in the diplomatic French
% layer. Critical readings and proofs are isolated in the worker ledgers.
% PDF091 / logical printed page 62 is confirmed blank topology.

\providecommand{\GalSourcePageStart}[4]{}%
\providecommand{\GalSourcePageBoundary}[4]{\space}%
\providecommand{\GalSourceBlankPage}[4]{}%
\providecommand{\GalPrintedError}[2]{#1}%
\providecommand{\GalSourceError}[2]{#1}%
\providecommand{\GalWitnessVariant}[2]{#1}%
\providecommand{\GalNotationAnomaly}[2]{#1}%
\providecommand{\GalHistoricalFootnoteMark}[1]{\textsuperscript{(#1)}}%
\providecommand{\GalHistoricalFootnoteText}[2]{\footnotetext{#2}}%

\GalSourcePageStart{080}{L0079}{51}{51}


% ENSEG:W10:PDF080:0001
\begin{center}
{\scriptsize\bfseries ON PRIMITIVE EQUATIONS\par}
\vspace{0.45em}
{\scriptsize\bfseries THAT ARE SOLVABLE BY RADICALS \GalHistoricalFootnoteMark{*}\par}
\vspace{0.45em}
{\scriptsize\bfseries (Fragment.)\par}
\end{center}
% ENSEG:W10:PDF080:0002
\GalHistoricalFootnoteText{*}{\emph{See} Note 1 on page 33.}

% ENSEG:W10:PDF080:0003
Let us seek, in general, the circumstances in which a primitive equation is
solvable by radicals. We can at once establish a general characteristic based
on the degree of such equations. It is this: {\itshape For a primitive equation
to be solvable by radicals, its degree must have the form $p^{\nu}$, where $p$
is prime.} It follows immediately that when an irreducible equation whose
degree has unequal prime factors is to be solved by radicals, this can be done
only by the decomposition method due to M.~Gauss; otherwise the equation will
be insoluble.

% ENSEG:W10:PDF080:0004
To establish the general property just stated for primitive equations solvable
by radicals, we may suppose that the equation to be solved is primitive but
ceases to be so upon adjoining a single radical. In other words, we may suppose
that, with $n$ prime, the group of the equation divides into $n$ conjugate
irreducible but nonprimitive groups. For, unless the degree of the equation is
prime, such a group will always occur in the sequence of decompositions.

% ENSEG:W10:PDF080:0005
Let $N$ be the degree of the equation, and suppose that after extracting a root
of prime degree $n$, it becomes nonprimitive and divides into $Q$ primitive
equations of degree $P$ by means of a single equation of degree $Q$.

% ENSEG:W10:PDF080:0006
If we call the group of the equation $G$, this group must divide into $n$
conjugate nonprimitive groups in which the letters are arranged in systems,
each consisting of $P$ connected letters. Let us see in how many ways this can
be done.

% ENSEG:W10:PDF080:0007
Let $H$ be one of the conjugate nonprimitive groups. It is easy to
\GalSourcePageBoundary{081}{L0080}{52}{52}
see that in this group any two letters chosen at will belong to some system of
$P$ connected letters, and belong to only one such system.

% ENSEG:W10:PDF081:0008
\GalWitnessVariant{For, in the first place}{W10-WV001}, if there were two letters that
could not belong to the same system
of $P$ connected letters, the group $G$---which is such that any one of its
substitutions transforms all the substitutions of $H$ into one another---would
be nonprimitive, contrary to the hypothesis.

\GalCriticalNote{W10-WV001}{Liouville's expanded transition}{The phrase
``For, in the first place'' follows the 1897 edition. Tannery reports only the
ordinal ``1\textsuperscript{o}'' in Galois's manuscript and attributes the
expanded connective to Liouville. The translation therefore keeps the printed
wording.}

% ENSEG:W10:PDF081:0009
Second, if two letters belonged to several different systems, it would follow
that the groups corresponding to the various systems of $P$ connected letters
were not primitive, again contrary to the hypothesis.

% ENSEG:W10:PDF081:0010
Accordingly, let
\[
\begin{array}{*{5}{l}}
a_{0}, & a_{1}, & a_{2}, & \dots, & a_{P-1}, \\
b_{0}, & b_{1}, & b_{2}, & \dots, & b_{P-1}, \\
c_{0}, & c_{1}, & c_{2}, & \dots, & c_{P-1}
\end{array}
\]
be the $N$ letters; suppose that each horizontal row represents a system of
connected letters. Let
\[
\GalNotationAnomaly{a_{0}}{W10-NA001},\quad a_{0,1}, \quad a_{0,2}, \ \dots, \quad a_{0,P-1}
\]
\GalCriticalNote{W10-NA001}{A possible missing second index}{The 1897 edition
prints $a_0$ before $a_{0,1},\ldots,a_{0,P-1}$. The pattern suggests
$a_{0,0}$, and a later control text silently expands it, but no independent
witness proves that a second index was omitted. The printed form is therefore
retained.}
be $P$ connected letters, all situated in the first vertical column. (It is
clear that we may arrange this by interchanging the order of the horizontal
rows.) Likewise, let
\[
a_{1,0}, \quad a_{1,1}, \quad a_{1,2}, \quad a_{1,3}, \ \dots, \quad a_{1,P-1}
\]
be $P$ connected letters, all situated in the second vertical column, so that
\[
a_{1,0}, \quad a_{1,1}, \quad a_{1,2}, \quad a_{1,3}, \ \dots, \quad a_{1,P-1}
\]
respectively belong to the same horizontal rows as
\[
a_{0,0}, \quad a_{0,1}, \quad a_{0,2}, \quad a_{0,3}, \ \dots, \quad a_{0,P-1};
\]
and similarly let the systems of connected letters be
\[
\begin{array}{*{6}{l}}
a_{2,0}, &a_{2,1}, &a_{2,2}, &a_{2,3}, & \dots, &a_{2,P-1}, \\
a_{3,0}, &a_{3,1}, &a_{3,2}, &a_{3,3}, & \dots, &a_{3,P-1}, \\
\dots,  & \dots,  & \dots, & \dots, & \dots, & \dotfill,
\end{array}
\]
\GalSourcePageBoundary{082}{L0081}{53}{53}
we thereby obtain $P^{2}$ letters in all. If the total number of letters is not
exhausted, we introduce a third index, so that
\[
a_{m,n,0}, \quad a_{m,n,1}, \quad a_{m,n,2}, \quad a_{m,n,3},\ \dots, \quad a_{m,n,P-1}
\]
is in general a system of connected letters. We thereby reach the conclusion
that $N=P^{\mu}$, where $\mu$ is the number of distinct indices required. The
general form of the letters will be
\[
\GalPrintedError{a_{k_{1}, k_{1}, k_{3}, \dots, k_{\mu}}}{W10-PE001},
\]
where $k_{1},k_{2},k_{3},\dots,k_{\mu}$ are indices, each of which may take the
$P$ values $0,1,2,3,\dots,P-1$.

% ENSEG:W10:PDF082:0011
The manner in which we proceeded also shows that every substitution in the
group $H$ has the form
\[
\bigl[a_{k_{1}, k_{2}, k_{3} \dots k_{\mu}}, \quad
      \GalPrintedError{a_{\phi(k_{1}), \psi(k)_{2}, \chi(k_{3}) \dots, \sigma(k_{\mu})}}{W10-PE002}
\bigr],
\]
because each index corresponds to a system of connected letters.

% ENSEG:W10:PDF082:0012
If $P$ is not prime, we reason about the permutation group of any one of the
systems of connected letters as we did about $G$, replacing each index by a
certain number of new indices. We find $P=R^{\alpha}$, and so on; ultimately
$N=p^{\nu}$, where $p$ is prime.

\GalCriticalNote{W10-DA001}{Why the degree is a prime power}{Choose a minimal
nontrivial normal subgroup $N$. Solvability makes $N$ elementary abelian; its
orbits are blocks, so primitivity makes it transitive. Because it is abelian
and the action is faithful, $N$ is regular. The point stabilizer acts
faithfully and irreducibly, whence the degree is $p^d$ and the group is affine.
This supplies the missing argument under the inherited finite, solvable,
primitive hypotheses.}

\GalCriticalNote{W10-PE001}{The second index is $k_{2}$}{The second printed
$k_{1}$ must be $k_{2}$, as required by the index declaration and confirmed by
Tannery's 1908 collation. With this correction, the coordinate argument used
in the prime-power theorem is restored.}
\GalCriticalNote{W10-PE002}{Correct placement of $\psi$}{Coordinate
parallelism and Tannery's collation establish $\psi(k_{2})$ in place of the
printed $\psi(k)_{2}$. This restores the coordinate argument used in the
prime-power theorem.}

\GalCriticalNote{W10-WV002}{Corrections recorded by Tannery}{Tannery's 1908
collation explicitly records the four nearby corrections to the indices,
coordinate pair, and affine formula. This supplies their transmission history
but does not authorize silently changing the translated 1897 text.}


\begin{center}\bfseries Primitive equations of degree $p^{2}$.\end{center}

% ENSEG:W10:PDF082:0013
Let us pause to treat at once the primitive equations of degree $p^{2}$, where
$p$ is an odd prime. (The case $p=2$ has been examined.) If an equation of
degree $p^{2}$ is solvable by radicals, suppose first that it becomes
nonprimitive upon extracting a radical.

% ENSEG:W10:PDF082:0014
Let $G$, then, be a primitive group on $p^{2}$ letters which divides into $n$
nonprimitive groups conjugate to $H$.

% ENSEG:W10:PDF082:0015
The letters in $H$ must necessarily be arranged
\GalSourcePageBoundary{083}{L0082}{54}{54}
as follows:
\[
\begin{array}{*{6}{l}}
  a_{0,0}, & a_{0,1}, & a_{0,2}, & a_{0,3}, & \dots, & a_{0,p-1},\\
  a_{1,0}, & a_{1,1}, & a_{1,2}, & a_{1,3}, & \dots, & a_{1,p-1},\\
  a_{2,0}, & a_{2,1}, & a_{2,2}, & a_{2,3}, & \dots, & a_{2,p-1},\\
  \dots,  & \dots,  & \dots,  & \dots,  & \dots, & \dotfill, \\
  a_{p-1,0}, & a_{p-1,1}, & a_{p-1,2}, & a_{p-1,3}, & \dots, & \rlap{$a_{p-1,p-1},$}
\end{array}
\]
each horizontal row and each vertical column being a system of connected letters.

% ENSEG:W10:PDF083:0016
If the horizontal rows are permuted among themselves, the group obtained,
being primitive and of prime degree, must contain only substitutions of the form
\[
\GalPrintedError{(a_{k_{1},k_{2}},\ a_{mk_{1}+n k_{2}})}{W10-PE003},
\]
the indices being taken modulo $p$.

The same will hold for the rows, which can yield only substitutions of the form
\[
(a_{k_{1},k_{2}},\ a_{k_{1}, qk_{2}+r}).
\]
Thus, finally, every substitution in $H$ has the form
\[
\GalPrintedError{(a_{k_{2},k_{3}},\ a_{m_{1} k_{1}+n_{1}, m_{2} k_{2}+n_{2}})}{W10-PE004}.
\]
If a group $G$ divides into $n$ groups conjugate to the one just described,
every substitution in $G$ must transform into one another the cyclic
substitutions of $H$, all written as follows:
\[
\tag{a}
(a_{k_{1},k_{2}},\ \dots,\ a_{k_{1}+\alpha_{1}, k_{2}+\alpha_{2}},\ \dots).
\]
Suppose, then, that one substitution of $G$ is formed by replacing, respectively,
\begin{align*}
k_{1} & \quad\text{by}\quad \phi_{1} (k_{1}, k_{2}),\\
k_{2} & \quad\text{by}\quad \phi_{2} (k_{1}, k_{2}).
\end{align*}
If in $\phi_{1},\phi_{2}$ we substitute $k_{1}+\alpha_{1}$ and
$k_{2}+\alpha_{2}$ for $k_{1},k_{2}$, the results must have the form
\[
\phi_{1} + \beta_{1},\quad \phi_{2} + \beta_{2},
\]
\GalSourcePageBoundary{084}{L0083}{55}{55}
and from this it is easy to conclude immediately that all substitutions of $G$
must be included in the formula
\[
\tag{A}
\GalPrintedError{(a_{k_{1}, k_{2}},\ a_{m_{1} k_{1} + n_{1} k + \alpha_{1} m_{2} k_{1} + n_{2} k_{2} + \alpha_{2}})}{W10-PE005},
\]
\GalCriticalNote{W10-PE003}{Restoring the second coordinate}{The second
coordinate has been dropped. The correct pair is
$(a_{k_{1},k_{2}},a_{mk_{1}+n,k_{2}})$, as required by the two-coordinate
construction and confirmed by Tannery. This repair restores the affine action
and the later counting argument.}
\GalCriticalNote{W10-PE004}{The correct coordinate pair}{The established
domain uses $(k_{1},k_{2})$; the printed $(k_{2},k_{3})$ does not belong to it.
The local two-index construction therefore forces $(k_{1},k_{2})$. No
independent historical correction was found, but this repair restores the
affine action.}
\GalCriticalNote{W10-PE005}{The correct affine formula}{The two-coordinate
image is
$a_{m_{1}k_{1}+n_{1}k_{2}+\alpha_{1},\,m_{2}k_{1}+n_{2}k_{2}+\alpha_{2}}$.
It is the unique form compatible with the finite-difference condition and is
confirmed by Tannery's collation. Under the inherited hypotheses established
above, the affine form and its group count survive.}
\GalCriticalNote{W10-TERM001}{What ``order'' means here}{Here ``order''
translates the 1897 use of \emph{ordre} for support size: the number of letters
moved, not the group-theoretic order of a substitution. Thus two, one, or no
fixed projective letters correspond to $p-1$, $p$, or $p+1$ moved letters.}
Now we know, from no.~\GalHistoricalFootnoteMark{*}, that the substitutions of
$G$ can involve only $p^{2}-1$ or $p^{2}-p$ letters. It cannot be $p^{2}-p$,
since in that case $G$ would be nonprimitive. Thus, if in $G$ we consider only
the permutations in which, for example, the letter $a_{0,0}$ always retains
the same place, we have only substitutions of order $p^{2}-1$ among the other
$p^{2}-1$ letters.
\GalHistoricalFootnoteText{*}{This memoir follows a work of Galois that I do
not possess; it is therefore impossible for me to identify the memoir cited
here and below. \textup{(A.~Ch.)}}

% ENSEG:W10:PDF084:0017
But let us recall that it was merely for purposes of the proof that we supposed
the primitive group $G$ to divide into conjugate nonprimitive groups. Since
this condition is by no means necessary, the groups will often be much more
composite.

We must therefore determine when these groups can admit substitutions in which
only $p^{2}-p$ letters vary, and this investigation will occupy us for some time.

Let $G$ be a group containing some substitution of order $p^{2}-p$. I first
assert that every substitution in this group is linear, that is, of form (A).

This is known to be true for substitutions of order $p^{2}-1$; it remains only
to prove it for those of order $p^{2}-p$. Consider, then, a group in which all
the substitutions are $m$ of order $p^{2}$ or order $p^{2}-p$. (\emph{See}
the passage cited.)

The $p$ letters that remain fixed in a substitution of order $p^{2}-p$ must
then be connected letters.

Suppose these connected letters are
\[
a_{0,0},\quad a_{0,1},\quad a_{0,2},\quad \dots, \quad a_{0,p-1}.
\]
We may derive all substitutions in which these $p$ letters do not change place
from substitutions
\GalSourcePageBoundary{085}{L0084}{56}{56}
of the form
\[
(a_{k_{1},k_{2}},\ a_{k_{1},\phi k_{2}})
\]
and from substitutions of order $p^{2}-p$ whose period has $p$ terms.
(\emph{See} again the passage cited.)

For the group to have the desired property, the former must necessarily reduce
to the form
\[
(a_{k_{1},k_{2}},\ a_{k_{1},m k_{2}})
\]
by what was seen for equations of degree $p$.

As for substitutions whose period has $p$ terms, since they are conjugate to
the preceding ones, we may suppose a group containing them but not the former.
They must therefore transform the cyclic substitutions (a) into one another,
and hence they too are linear.

\GalCriticalNote{W10-DA002}{The affine form needs solvability}{Under the
inherited finite, solvable, primitive hypotheses, the group has a regular
elementary-abelian socle, so every element acts as $x\mapsto Ax+b$.
Primitivity alone is not enough: the natural action of $S_{p^2}$ for
$p^2\ge5$ is primitive but not affine. Formula (A) therefore survives only in
the inherited solvable setting and with the corrected affine formula.}

We have thus reached the conclusion that the primitive permutation group on
$p^{2}$ letters must contain only substitutions of form (A).

% ENSEG:W10:PDF085:0018
Now take the full group obtained by applying to the expression
\[
a_{k_{1},k_{2}}
\]
all possible linear substitutions, and seek the subgroups of this group that
can have the property required for solvability of equations.

What, first, is the total number of linear substitutions? It is clear that not
every transformation of the form
\[
k_{1},\ k_{2},\quad m_{1}k_{1} + n_{1}k_{2} + \alpha_{1},\ m_{2}k_{1} + n_{2}k_{2} + \alpha_{2}
\]
is thereby a substitution; in a substitution, each letter of the first
permutation must correspond to exactly one letter of the second, and conversely.

Thus, if we take any letter $a_{l_{1},l_{2}}$ of the second permutation and
trace it back to the corresponding letter in the first, we must find a letter
$a_{k_{1},k_{2}}$ whose indices $k_{1},k_{2}$ are completely determined.
Hence, for arbitrary $l_{1},l_{2}$, the two equations
\[
m_{1}k_{1} + n_{1}k_{2} + \alpha_{1} = l_{1}, \quad
m_{2}k_{1} + n_{2}k_{2} + \alpha_{2} = l_{2},
\]
\GalSourcePageBoundary{086}{L0085}{57}{57}
must yield finite and determinate values of $k_{1},k_{2}$. The condition for
such a transformation actually to be a substitution
\GalWitnessVariant{is}{W10-WV003} therefore that
$m_{1}n_{2}-m_{2}n_{1}$ be neither zero nor divisible by the modulus $p$,
which is the same thing.

% ENSEG:W10:PDF086:0019
I now assert that although, as we shall see, this linear substitution group
does not always belong to equations solvable by radicals, it nevertheless has
the following property: if any one of its substitutions has $n$ fixed letters,
then $n$ divides the number of letters. Indeed, whatever the number of letters
remaining fixed, this condition can be expressed by linear equations giving
all the indices of one fixed letter in terms of some number of them. Giving
each index that remains arbitrary its $p$ values yields $p^{m}$ systems of
values, where $m$ is some number. \GalSourceError{In the case before us,
$m$ is necessarily $<2$ and is consequently $0$ or $1$.}{W10-SE001} Therefore
the number of substitutions cannot exceed
\[
p^{2}(p^{2} - 1)(p^{2} - p).
\]

\GalCriticalNote{W10-SE001}{Exclude the identity map}{For a nonidentity affine
map of $\mathbb F_{p}^{2}$, the fixed set is empty or is an affine subspace of
dimension $0$ or $1$; the identity has fixed dimension $2$. The printed bound
therefore needs the word ``nonidentity.'' The divisibility conclusion for the
number of fixed letters remains valid, including for the identity.}

\GalCriticalNote{W10-WV003}{Chevalier supplied ``is''}{The word ``is'' appears
in the 1897 sentence. Tannery reports that it was absent from Galois's
manuscript and supplied by Chevalier. The translation retains the printed
syntax.}

\GalCriticalNote{W10-WV004}{Chevalier's expanded wording}{According to
Tannery, the 1897 phrase ``is necessarily $<2$ and is consequently $0$ or
$1$'' expands Galois's terser wording. The translation follows the printed
edition rather than combining it with the manuscript sentence.}

% ENSEG:W10:PDF086:0020
Now consider only the linear substitutions in which $a_{0,0}$ does not move.
If in this case we find the total number of permutations in the group
containing all possible linear substitutions, it will suffice to multiply
this number by $p^{2}$.

First, by substituting $p$ for the index $k_{2}$, all substitutions of the form
\[
(a_{k_{1}, k_{2}},\ a_{m_{1}k_{1}, k_{2}})
\]
give $p-1$ substitutions in all. We obtain $p^{2}-p$ of them by adding the
term $m_{2}k_{1}$ to $k_{2}$, as follows:
\[
\tag{m'}
(k_{1},\ k_{2},\ m_{1}k_{1},\ m_{2}k_{1} + k_{2}).
\]

On the other hand, it is easy to find a linear group of $p^{2}-1$ permutations
such that in each substitution all letters except $a_{0,0}$ move. For, replacing
the double index $k_{1},k_{2}$ by the single index $k_{1}+ik_{2}$,
\GalSourcePageBoundary{087}{L0086}{58}{58}
where $i$ is a primitive root of
\[
x^{p^{2}-1}\! - 1 = 0\ (\mod p),
\]
it is clear that every substitution of the form
\[
[a_{k_{1}+k_{2} i},\ a_{(m_{1}+m_{2}i)(k_{1}+k_{2}i)}]
\]
is linear; \GalSourceError{but in these substitutions no letter remains in the
same place, and there are $p^{2}-1$ of them.}{W10-SE002}

\GalCriticalNote{W10-SE002}{Counting nonidentity scalars}{The scalar subgroup
has $p^{2}-1$ elements, including the identity. Exactly $p^{2}-2$ nonidentity
scalars move every nonzero letter, while every scalar fixes $a_{0,0}$. The
subgroup order and the later count of the full affine group remain unchanged;
only the universal claim about moved letters must exclude the identity.}

% ENSEG:W10:PDF087:0021
We therefore have a system of $p^{2}-1$ permutations such that in each
substitution all letters except $a_{0,0}$ move. Combining these substitutions
with the $p^{2}-p$ mentioned above gives
\[
(p^{2} - 1)(p^{2} - p) \text{ substitutions}.
\]

Now we saw \textit{a priori} that the number of substitutions in which
$a_{0,0}$ remains fixed could not exceed $(p^{2}-1)(p^{2}-p)$. It is therefore
exactly $(p^{2}-1)(p^{2}-p)$, and the full linear group has altogether
\[
p^{2}(p^{2} - 1)(p^{2} - p) \text{ permutations}.
\]

% ENSEG:W10:PDF087:0022
It remains to find the subgroups of this group having the property required for
solvability by radicals. For this purpose we shall make a transformation aimed
at lowering as far as possible the general equations of degree $p^{2}$ whose
group is linear.

First, since the cyclic substitutions of such a group are transformed into
one another by every other substitution of the group, we may lower the degree
of the equation by one and consider an equation of degree $p^{2}-1$ whose
group has only substitutions of the form
\[
(b_{k_{1}, k_{2}},\
 \GalWitnessVariant{b_{m_{1}k_{1}+n_{1}k_{2},\: m_{2}k_{1}+n_{2}k_{2}}}{W10-WV005}),
\]
the $p^{2}-1$ letters being
\[
\begin{array}{*{5}{l}}
         &b_{0,1}, &b_{0,2}, &b_{0,3}, &\dots,\\
b_{1,0}, &b_{1,1}, &b_{1,2}, &b_{1,3}, &\dots,\\
b_{2,0}, &b_{2,1}, &b_{2,2}, &b_{2,3}, &\dots,\\
\dots,  &\dots,  &\dots,  &\dots,  &\dots.
\end{array}
\]

\GalCriticalNote{W10-WV005}{The 1897 formula is already correct}{Tannery notes
that this two-coordinate formula in the 1897 edition corrects an error in the
Chevalier--Liouville transmission. The printed 1897 formula should therefore
be retained; it needs no editorial replacement.}
\GalSourcePageBoundary{088}{L0087}{59}{59}

% ENSEG:W10:PDF088:0023
I now observe that this group is nonprimitive, so all letters for which the
ratio of the two indices is the same are connected letters. Replacing each
system of connected letters by a single letter gives a group whose
substitutions all have the form
\[
\biggl(b_{\frac{k_{1}}{k_{2}}},\
      b_{\frac{m_{1} k_{1} + n_{1} k_{2}}{m_{2} k_{1} + n_{2} k_{2}}}\biggr),
\]
where $\dfrac{k_{1}}{k_{2}}$ are the new indices. Replacing this ratio by one
index $k$, we see that the $p+1$ letters are
\[
b_{0}, \quad b_{1}, \quad b_{2}, \quad b_{3}, \quad \dots, \quad b_{p-1}, \quad b_{\frac{1}{0}},
\]
and the substitutions have the form
\[
\biggl(k,\ \frac{mk + n}{rk + s}\biggr).
\]

Let us determine how many letters remain fixed in each of these substitutions;
for this we must solve the equation
\[
\GalPrintedError{(rk + s)k - m(mk + n) = 0}{W10-PE006},
\]
which has two, one, or no roots according as $(m-s)^{2}+4nr$ is a quadratic
residue, zero, or a quadratic nonresidue. In these three cases respectively,
the substitution has order $p-1$, $p$, or $p+1$.

\GalCriticalNote{W10-PE006}{Removing an extra factor}{The fixed-point equation
is $(rk+s)k-(mk+n)=0$. The stated discriminant is the discriminant of this
polynomial, not of the printed polynomial with an extra factor $m$. The same
extra factor already appears in the 1846 publication, which shows that the
error was inherited but does not make it correct. Removing it restores the
discriminant and the three cases for the support; the later $p=3$ inference
remains independently invalid.}

\GalCriticalNote{W10-WV006}{The extra factor was inherited from 1846}{The
extra factor $m$ in the fixed-point equation already appears in the 1846
publication. This establishes how the error reached the 1897 edition, but does
not validate the formula; the corrected equation is given in the preceding
note.}

As a type for the first two cases we may take substitutions of the form
\[
(k,\ mk + n),
\]
in which only the letter $b_{\frac{1}{0}}$ does not move. It follows that the
total number of substitutions in the reduced group is
\[
(p + 1)p(p - 1).
\]

% ENSEG:W10:PDF088:0024
Having reduced this group in this way, we shall now treat it generally. We
shall first determine when a subgroup containing substitutions of order $p$
can belong to an equation solvable by radicals.

In this case the equation would be primitive and could not be
\GalSourcePageBoundary{089}{L0088}{60}{60}
solvable by radicals unless $p+1=2^{n}$ for some $n$.

We may suppose that the group contains only substitutions of order $p$ and
order $p+1$. Consequently all substitutions of order $p+1$ are similar and
have a period of two terms.

Take the expression
\[
\biggl(k,\ \frac{mk + n}{rk + s}\biggr),
\]
and determine when this substitution can have a period of two terms. For this,
the inverse substitution must coincide with it. The inverse substitution is
\[
\biggl(k,\ \frac{-sk + n}{rk - m}\biggr).
\]

Thus $m=-s$, and all substitutions in question are
\[
\biggl(k,\ \frac{mk + n}{k - m}\biggr),
\]
or equivalently
\[
k,\ m + \frac{N}{k - m},
\]
where $N$ is some number, the same for every substitution, since these
substitutions must be transformed into one another by all substitutions of
order $p$, $(k,k+m)$. Moreover, these substitutions must be conjugate to one
another. Thus, if
\[
\biggl(k,\ m + \frac{N}{k - m}\biggr), \qquad
\biggl(k,\ n + \frac{N}{k - n}\biggr)
\]
are two such substitutions, we must have
\[
n + \frac{N}{\dfrac{N}{k-m} + m - n} =
m + \frac{N}{\dfrac{N}{k-n} + n - m},
\]
namely,
\[
(m - n)^{2} = 2N.
\]

\GalCriticalNote{W10-DA003}{Conjugate substitutions need not commute}{The
nineteenth-century phrase ``conjoined substitutions'' does not imply pairwise
commutation, nor does modern conjugacy. On $\mathbf P^1(\mathbf F_7)$, the maps
$T_m(k)=m+1/(k-m)$ give translation-conjugate transformations $T_0$ and $T_1$
that do not commute. The displayed pairwise equality is therefore unsupported,
and the final inference $p=3$ does not follow.}

Thus the difference between two values of $m$ can assume
\GalSourcePageBoundary{090}{L0089}{61}{61}
only two distinct values; hence $m$ can have no more than three values, and
therefore $p=3$. Thus only in this case can the reduced group contain
substitutions of order $p$.

Indeed, the reduced equation is then of the fourth degree and is consequently
solvable by radicals.

We therefore know that, in general, the substitutions of our reduced group
must include none of order $p$. Can there be any of order $p-1$? This is what
I shall investigate\GalHistoricalFootnoteMark{1}.

\GalHistoricalFootnoteText{1}{I searched Galois's papers in vain for the
continuation of what has just been read. \textup{(A.~Ch.)}}

\GalCriticalNote{W10-DA004}{What survives after the corrections}{The corrected
indices and coordinate pairs restore the coordinate maps; the corrected
affine formula is valid under the inherited finite, solvable, primitive
hypotheses. Excluding the identity preserves the fixed-letter divisibility
result, and counting only nonidentity scalars changes that count to $p^2-2$
without changing either the subgroup order or the full affine-group order.
Removing the extra factor restores the fixed-point equation and the three
support cases. Only the final inference $p=3$, which rests on the unsupported
commutativity premise, remains blocked.}

\par\vspace{1.0em}\noindent\dotfill\par
\vspace{0.4em}\hrule\vspace{2.6em}
\begin{center}\bfseries END.\end{center}

\GalSourceBlankPage{091}{L0090}{62}{not visibly printed}
